% Book-local P07.1 positive case, 20 September 2026.
% This isolated statement makes no novelty or priority claim.

\begin{proposition}
\label{prop:p071-trivalent-affine-cross-collar-smoothing}
Let $E=F=\mathbf R^2$ and put
\[
 v_1=w_1=(1,0),\qquad
 v_2=w_2=(0,1),\qquad
 v_3=w_3=(-1,-1).
\]
Let $S$ and $T$ be the corresponding open exact trivalent stars, with
$p_i(r)=rv_i$ and $q_j(s)=sw_j$ for $0\leq r,s<R$.  On $S\times T$ let
$\mathcal G_0$ be the local generator class whose members are real-valued,
have smooth jointly convex restrictions to the nine branch rectangles,
have common smooth restrictions on the two families of node axes, and
satisfy
\begin{align}
 \sum_{i=1}^3\partial_rg_{ij}(0,s)&=0
       &&(j=1,2,3),
 \label{eq:p071-cross-collar-first-balance}\\
 \sum_{j=1}^3\partial_sg_{ij}(r,0)&=0
       &&(i=1,2,3).
 \label{eq:p071-cross-collar-second-balance}
\end{align}

Fix $0<R_0<R$ and $0<\delta<R_0$.  Suppose that $u$ is a continuous
real-valued function on the closed product substar with $0\leq r,s\leq
R_0$, that every restriction
\[
                         u_{ij}(r,s)=u(p_i(r),q_j(s))
\]
is convex and piecewise affine on $[0,R_0]^2$, and that there are
$c\in\mathbf R$, $A\in E^*$, and $B\in F^*$ such that
\begin{equation}
 u_{ij}(r,s)=c+A(rv_i)+B(sw_j)
 \quad\text{whenever}\quad \min\{r,s\}<\delta.
 \label{eq:p071-cross-collar-affine-data}
\end{equation}
Then on the fixed open product substar
\[
 V=\{p_i(r):0\leq r<R_0\}\times
   \{q_j(s):0\leq s<R_0\}
\]
there are $G_n\in\mathcal G_0(V)$ such that
\begin{equation}
                         G_{n+1}\leq G_n,
             \qquad G_n\longrightarrow u
 \label{eq:p071-cross-collar-monotone-approximation}
\end{equation}
uniformly on $V$.  Each $G_n$ agrees with one ambient affine function on
the same cross-collar $\min\{r,s\}<\delta$.

Consequently every such $u$ belongs already to the first decreasing-limit
stage of the generated two-star sheaf.  The proposition permits arbitrary
finite convex polyhedral coupling away from the singular axes, but it does
not remove the affine cross-collar hypothesis and does not retire P07.1.
\end{proposition}
